\section{Level 3: Action-Conditioned Simulation}
\label{sec:level3}

\subsection{Capability Definition and Adjacent Boundary}

Level~3 models a future state, observation, or outcome after an explicit action has been set. Operationally, the model represents $p(s_{t+1},y_{t+1}\mid s_t,A_t{=}a)$, where
$A_t$ is a selectable clinical, acquisition, surgical, physiological, or management action.. This is the minimum contract for an action-aware medical world model, where the settable action enters the transition, and changes an evaluated future.

Three conditions define the boundary from Level~2. First, the action must be settable at inference for the same pre-action state or history. Second, it must enter the transition rather than serve only as a caption, classifier target, reward, or narrative prompt. Third, the system must predict a meaningful post-action future. Time, disease stage, desired output sketches, action-class text, and recorded clinician behavior do not pass the test on wording alone. Acquisition actions qualify when they control the future observation process, even if the patient state itself is unchanged. SATO-V records that narrower scope.

The upper boundary is Level~4 counterfactual outcome reasoning. Level~3 permits descriptive intervention-conditioned simulation. A treatment token may change a generated trajectory because treatment assignment is associated with disease severity and clinician behavior. Level~4 requires an explicit reason why outcomes under alternative treatments can be compared for the same patient or state. This separation allows action-conditioned systems to receive appropriate credit without importing causal language that their data and evaluation do not support.

\subsection{Clinical Evidence Landscape}

Treatment-aware imaging and disease-progression models are prominent Level~3 examples. A treatment label, treatment day, dose, or protocol variable conditions a diffusion or latent progression model that generates a post-treatment tumor or anatomical state \citep{yang2025medical,liu2019novel}. These models move beyond natural-history forecasting because an intervention enters the transition. Their clinical interpretation depends on the granularity and provenance of that intervention. A coarse label can implement the task form while remaining a proxy for patient subgroup, disease stage, or local practice.

The action-conditioned imaging cluster also includes predictive sensing in chest radiography, glioblastoma world models, multimodal physiological intervention simulation, and ontology-guided disease-risk twins \citep{yang2025xray2xray,yang2026x,wang2026brain,lutsker2026simulating,li2026modeling}. These systems expose different actions: acquisition geometry, treatment context, physiological intervention, or public-health scenario. Grouping them at Level~3 does not assert that the actions have equal clinical meaning. It records the narrower common contract that changing a settable input changes an evaluated future state or observation.

Imaging acquisition provides a observation-level action, where rotation, view selection, or probe motion changes the future measurement rather than the underlying disease process. Xray2Xray uses a selectable acquisition rotation to change the predicted radiographic observation \citep{yang2025xray2xray}. EchoWorld conditions future echocardiographic views on probe motion \citep{yue2025echoworld}. These actions are executable and directly connected to the observation process. They do not necessarily change patient biology, but they can support acquisition guidance, training, and embodied perception. This distinction shows why application and action semantics must be recorded alongside the capability level.

Surgery and robotics supply similarly explicit environmental actions. SWoMo models surgical tool-motion increments in a neuro-symbolic simulator, while surgical video world models condition future scenes on tool actions \citep{sivakumar2026swomo,rapuri2026saw,lin2024world}. The action can be instrumented, repeated, and evaluated at high temporal resolution. Their evaluation should therefore extend from visual continuity to geometry, tissue response, force, task success, and safety-relevant behavior.

The embodied cluster is broader than controllable video. Surgical vision and unified surgical world models combine scene understanding with future generation; echocardiography systems use probe motion or structural registration as the action; real-time soft-tissue simulation and long-horizon surgical modeling target physical response; large embodiment datasets support model training rather than establishing a capability by themselves \citep{koju2025surgical,unknown2025UnifiedSurgicalWorld,jiang2024cardiac,greifeneder2026autoregressive,jiang2024structure,nelson2026open,kang2026dreamreg,pan2026surgvista}. EchoWorld provides a particularly clear acquisition-action example because probe motion is explicitly coupled to the predicted view \citep{yue2025echoworld}. Across this cluster, evaluation must move from frame realism to geometry, tissue response, registration accuracy, and task-relevant safety.

EHR and physiological models occupy a more ambiguous part of Level~3. A patient simulator can qualify when a medication, dose, fluid, vasopressor setting, or management choice is deliberately supplied and the model predicts subsequent physiology. By contrast, an autoregressive model that emits the next medication or procedure as part of the record remains Level~2. Physiological priors can strengthen action semantics when interventions have known mechanisms, but retrospective treatment data still reflect selection and support limitations. Embodiment and mechanism clarify the intervention; they do not complete the causal argument.

Recent patient-trajectory and twin systems illustrate this range. EHRWorld and ChronoMedicalWorld expose longitudinal care dynamics; sepsis twins and Clin-JEPA-style patient representations connect structured history to downstream intervention modeling; cloud and standard-of-care twins place the patient model inside a broader care interface \citep{mu2026ehrworld,wang2026chronomedicalworld,lal2020development,yang2026clin,liu2019novel,bhattacharya2025soc}. The key audit question is whether an external action is supplied to the transition, not whether medications appear somewhere in the generated record.

Mechanistic and physiological twins provide more explicit state variables but still require empirical action support. Oncology twins with quantified uncertainty, stimulation-oriented virtual-brain twins, drug-aware physiological diffusion, latent physiological dynamics, and probabilistic cardiac twins all model intervention-sensitive futures at different temporal and biological scales \citep{pash2026predictive,wang2025virtual,shao2025mm,luo2026toward,giovanis2026probabilistic}. Their interpretability depends on which parameters are individualized and which intervention-response relationships are learned, assumed, or externally constrained.

Digital twins can appear at this level when they update a patient-specific state and simulate an intervention. The important distinction is implementation. A diagram of a future treatment interface is not action-conditioned evidence. A mechanistic or learned twin that accepts a defined regimen and produces a patient-specific rollout is. Patient updating, parameter uncertainty, and validation under observed interventions determine whether the twin is an operational simulator rather than an aspirational architecture.

\subsection{Representative Model Designs}

Diffusion, flow, and video generators often inject action through cross-attention, embeddings, or control channels. Their ability to model multimodal futures is attractive for treatment response and embodied scenes. The decisive design question is whether the action has a stable semantic interface and whether changing it alters outcome-relevant aspects of the rollout rather than superficial appearance. Action masking, swapping, and dose-response tests are therefore as important as image quality.

Autoregressive and recurrent patient models include actions in the event history or as separate inputs to a latent transition. Separating those roles is essential. If the model predicts an action and then feeds its own prediction back into the sequence, it may be simulating the care process. If an external action is set and the patient state is advanced conditionally, it implements the Level~3 contract. Latent state-space models can make the distinction clearer by separating observation, action, and transition variables, though identifiability and observability remain concerns.

Mechanistic, graph, and neuro-symbolic systems encode domain structure directly. Surgical simulators can impose geometry and tool constraints; physiological systems can use differential equations or known dose-response pathways; hybrid models can learn residual dynamics or patient-specific parameters. These constraints may improve extrapolation and make failures more interpretable. They can also create false confidence when parameters are weakly identified or when the mechanistic model omits clinically important heterogeneity.

The architecture should expose the action pathway. A reader should be able to identify where the action enters, what state variables it can affect, what actions are admissible, and how uncertainty propagates. An action embedding without this interface is difficult to distinguish from generic conditioning. For the same reason, architecture prevalence is less informative than whether each design makes the intervention contract inspectable.

\subsection{Evidence Quality and Failure Modes}

Level~3 validation should begin with action sensitivity. Holding the pre-action state fixed, the analysis should vary actions, compare against time-only and no-action baselines, mask or shuffle action inputs, and test whether changes occur in the expected variables and direction. Results should be stratified by common, rare, and unsupported actions. Dose-response structure, known physical constraints, and mechanistic anchors can provide additional checks. These analyses show whether action conditioning is operational; they still do not prove that the predicted difference is causal.

The characteristic overclaim is the move from ``conditioned on treatment'' to ``estimates treatment effect.'' In retrospective data, treatment choice encodes severity, biomarkers, contraindications, prior response, access, and clinician preference. A model can learn these associations and generate realistic treated trajectories while failing under a different policy. The same problem appears in surgery when action labels correlate with phase or scene, and in acquisition when a requested view can be generated without preserving three-dimensional consistency.

Other failure modes include weak action support, coarse or clinically implausible action spaces, and outcome mismatch. A model may accept actions rarely or never observed for the current state. A binary treatment label may hide dose, schedule, combination therapy, or timing. A generated image may be evaluated by perceptual similarity when the stated use requires tumor response or safety. Each limitation can coexist with a valid Level~3 task form, which is why capability and evidence quality are reported separately.

An action card would make these studies easier to assess. It should describe executability, timing, admissible range, provenance, empirical support, the state variables expected to respond, and the outcome used for validation. This information allows a descriptive simulator to be evaluated on its own terms and clarifies what further causal evidence would be needed.

\subsection{Level Takeaway and Next Threshold}

Level 3 marks the point where the modeled future becomes responsive to a selectable operation on the patient, measurement process, or clinical environment. The strongest examples arise where actions are instrumented or mechanistically defined. The main weakness arises when observational treatment labels are granted causal meaning without an identification argument. The appropriate response is not to exclude action-conditioned work, but to keep its claim at the level supported by the evidence.

\levelsummary
{Treatment, acquisition, surgical, and physiological actions can be integrated into learned or mechanistic transitions, creating useful action-responsive simulators.}
{Action conditioning is not treatment-effect estimation. In observational medicine, an explicit action input can still encode selection bias, practice patterns, and unsupported extrapolation.}
{To reach Level~4, a study must define the alternative-outcome estimand and support it through causal identification, structural assumptions, trial or quasi-experimental evidence, mechanistic constraints, or sensitivity analysis.}
